% page 523 du fac-simile (image p-522 du scan)
% bloc 170.2 x 237.7 mm ; marge gauche 31.8 mm
\pgc{10.160mm}{0.000mm}{
\l{\cel{57}515}
\l{}
\l{\sou{seciono}. (\sou{geom.}) Ligno, surfaco, segun qua su partigas du sur-}
\l{\cel{3}faci, du solidi, un solido ed un surfaco. -\cel{1}\sou{Secioni konala}:}
\l{\cel{3}secioni plana di kono rekta kun bazo cirkla (cirklo, elipso,}
\l{\cel{3}parabolo, hiperbolo). - DEFIRS.}
\l{}
\l{\sou{sedo}. (\sou{anat.}) La parto karnoza qua esas infra-rene, che la homo,}
\l{\cel{3}- dope che la animali. -}
\l{}
\l{\sou{sedentaria}. I. Qua sidas maxim-multa-tempe. - II. Qua chanjas nur}
\l{\cel{3}rare sua rezideyo. - DEFIS.}
\l{}
\l{\sou{sediciar}. (\sou{netrans}.)Rezistar publike la autoritatozi establisita.}
\l{\cel{3}- EFIS.}
\l{}
\l{\sou{sedimento}. I. Depozajo di la materii solida qua esis suspendita}
\l{\cel{3}en liquido. - II. (\sou{geol.}) Strato quan la aqui depozis lor sua}
\l{\cel{3}departo.- DEFIS.}
\l{}
\l{\sou{seduktar}. (\sou{trans., per)}. Venigar (ulu) ad su per charmar (lu).-}
\l{\cel{3}EFIS.}
\l{}
\l{\sou{segar}. (\sou{trans.}) Tranchar, o nur sekar, per instrumento qua konsis-}
\l{\cel{3}tas ek lamo metala, dina, kun aristo dentoza, muntita an mancho,}
\l{\cel{3}e qua, per movo alternanta o rotacala, entamas, dividas pokope}
\l{\cel{3}la materii harda (ligno, quadro, petro, e c.) - DI.}
\l{}
\l{\sou{seglo}. (\sou{navig}.)Peco de telo forta, qua konsistas (ordinare) ye}
\l{\cel{3}plura peci asemblita, quan onu ligas an la yardi di la masti,}
\l{\cel{3}e qua, desfaldite ed orientizite segun la vigoro e la direciono}
\l{\cel{3}di vento, recevas de olca impulso qua igas la navo avancar ta-}
\l{\cel{3}direcione o ca-direcione. - DE.}
\l{}
\l{\sou{segmento}. Parto determinita di objekto longa, lineatra, quan onu}
\l{\cel{3}dividas, longesale : segmento di voyo, di fervoyo (inter du}
\l{\cel{3}stacioni), di rivero, di kanalo (inter du sluzi), di filo tele-}
\l{\cel{3}grafala o telefonala (inter du stacioni o du posteni). - DEFIRS.}
\l{}
\l{\sou{seido}. Ta qua devotesas fanatike ad ulu, qua esas pronta agar o}
\l{\cel{3}facar irgo por obediar lu, mem kriminar. DFIS.}
\l{}
\l{\sou{segregacar}. (\sou{trans}.) Separar (parto) de la bloko, de la socio. -}
\l{\cel{3}DEFIS.- (antonimo :\cel{1}\sou{agregar}).-}
\l{}
\l{\sou{segun}. I. Egardante (ulo quo dependas de la personi, de la cir-}
\l{\cel{3}konstanci). II. Fidante su a (la arbitro da, la decido da,}
\l{\cel{3}la selekto da, la dico da). - IS.}
\l{}
\l{\sou{sejornar}. (\sou{netrans., en, che)}\cel{1}. Esar dum kelka tempo, ma ne multa}
\l{\cel{3}(en ta o ca loko, che ta o ca persono). - EF.}
\l{}
\l{\sou{sekar}.\cel{2}(\sou{trans.}) (\sou{kirurg.}) Dividar, plu o min profunde,(korpo)}
\l{\cel{3}per korpo plu harda quaes as taliita tale ke lu havas aristo}
\l{\cel{3}tranchiva (skalpelo, e c.) - EFIS.}
}
